% Accelerator map, execute fabric, handshake, integer units.
% Requires tikz_styles.tex.
\begin{figure}[p]
\centering
\begin{tikzpicture}[x=1cm,y=1cm]
  \node[ctl] (exec) at (6.2,3.55) {S\_EXEC};
  \node[acc, minimum width=22mm] (alu) at (1.4,2.15) {int ALU\\1 cycle};
  \node[acc, minimum width=18mm] (mul) at (4.2,2.15) {mul\\5 cyc};
  \node[acc, minimum width=18mm] (div) at (6.6,2.15) {div\\5--37};
  \node[acc, minimum width=18mm] (ipow) at (9.0,2.15) {int pow};
  \node[acc, minimum width=32mm] (fpu) at (11.8,2.15) {FPU};
  \node[acc, minimum width=28mm] (cont) at (1.8,0.45) {container FSM};
  \node[acc, minimum width=22mm] (sa) at (5.4,0.45) {STRACC};
  \node[acc, minimum width=22mm] (gc) at (8.4,0.45) {GC};
  \node[ex, minimum width=28mm] (excore) at (11.6,0.45) {excore};
  \draw[arr] (exec) -- (div);
  \draw[arr] (exec.west) -- ++(-0.35,0) |- (alu.north);
  \draw[arr] (exec.east) -- ++(0.45,0) |- (fpu.north);
  \draw[darr] (gc.east) -- node[anno, above] {before commit} (excore.west);
\end{tikzpicture}
\caption{Accelerators, together.
Three kinds of hardware sit off the simple pipeline.
The execute fabric is tag-routed arithmetic and stays inside \texttt{S\_EXEC}.
The container FSM, the string accelerator, and the collector are multi-cycle masters of the data port; the core is frozen in their state until they finish.
Excore is the firmware accelerator for the cases those units refuse to do in place: growing a container and the uncontaminated bulk updates.
Negative integer powers never reach the integer power unit; execute promotes both operands and reissues them on the FPU.}
\label{fig:accel-all}
\end{figure}

\begin{figure}[p]
\centering
\begin{tikzpicture}[node distance=3.2mm and 3.6mm]
  \node[rf] (rs) {rs1, rs2\\132-bit};
  \node[ctl, right=of rs] (td) {\texttt{tag\_decode}};
  \node[ctl, right=of td] (pr) {\texttt{promote}};
  \node[acc, below=of pr] (intu) {INT units};
  \node[acc, right=of intu] (fpu) {FPU};
  \node[hi, left=of intu] (bool) {BOOL};
  \node[mem, below=of intu] (str) {short-string\\fast path};
  \node[ctl, right=5mm of fpu] (pack) {pack\\tag $+$ value};
  \node[trap, below=of pack] (tr) {trap};
  \draw[arr] (rs) -- (td);
  \draw[arr] (td) -- (pr);
  \draw[arr] (pr) -- (intu);
  \draw[arr] (pr) -- (fpu);
  \draw[arr] (pr) -- (bool);
  \draw[arr] (td.south) -- ++(0,-0.9) -| (str.north);
  \draw[arr] (intu) -- (pack);
  \draw[arr] (fpu) -- (pack);
  \draw[arr] (bool.east) -- ++(0.3,0) |- (pack.west);
  \draw[arr] (str.east) -- ++(0.4,0) |- (pack.south);
  \draw[darr] (td.east) -- ++(0.25,0) |- (tr.north);
  \node[anno, anchor=west, align=left] at ($(pack.east)+(0.25,0.15)$)
    {BOOL$\to$INT\\BOOL$\to$FLOAT\\INT$\to$FLOAT};
\end{tikzpicture}
\caption{Execute fabric (\texttt{pycore\_exec}).
Tag decode chooses the unit, the promotion, and the result tag.
Promotion is a combinational rewrite of the 64-bit lane.
The integer ALU, BOOL operations, short-string compare and concatenation, and the FPU's compares and unary operations finish in the issue cycle.
Multiplier, divider, integer power, and the FPU hold \texttt{S\_EXEC} through \texttt{stall\_o}.
The result is written back as a fresh 132-bit entry.
An integer result is sign-extended from 64 bits into the high half.}
\label{fig:exec}
\end{figure}

\begin{figure}[htbp]
\centering
\begin{tikzpicture}[x=1.15cm,y=0.7cm]
  \draw[pyedge] (0,2.2) -- (10.2,2.2);
  \draw[pyedge] (0,1.1) -- (10.2,1.1);
  \draw[pyedge] (0,0) -- (10.2,0);
  \node[anchor=east, font=\tiny] at (-0.1,2.2) {start};
  \node[anchor=east, font=\tiny] at (-0.1,1.1) {busy};
  \node[anchor=east, font=\tiny] at (-0.1,0) {done};
  \draw[pyacc, line width=2.2pt] (0.6,2.2) -- (7.4,2.2);
  \draw[pyctl, line width=2.2pt] (0.8,1.1) -- (7.2,1.1);
  \draw[pyedge, line width=2.2pt] (7.0,0) -- (7.8,0);
  \foreach \x in {0.6,2.2,3.8,5.4,7.0} {
    \draw[pyedge, thin] (\x,-0.25) -- (\x,2.55);
  }
  \node[anno] at (3.8,-0.6) {level start \quad one-cycle done \quad stall $=$ start and not done};
  \node[anno] at (5.0,2.7) {dropping start aborts and returns the unit to idle};
\end{tikzpicture}
\caption{Shared handshake of every multi-cycle arithmetic unit.
\texttt{start} is a level.
The unit accepts it when idle and not already in the done cycle, so a requester that holds \texttt{start} and swaps operands the cycle after \texttt{done} gets back-to-back issue.
\texttt{busy} is the structural-hazard bit a later scoreboard would use.
Traps (\texttt{div\_zero}, integer overflow, FPU exception) replace \texttt{done} and are reported with \texttt{stall} low.
The core leaves \texttt{S\_EXEC} by dropping \texttt{start}, which is what aborts a unit today.}
\label{fig:hs}
\end{figure}

\begin{figure}[htbp]
\centering
\begin{tikzpicture}[node distance=3mm and 3mm]
  \node[rf] (a) {a[63:0]};
  \node[rf, right=of a] (b) {b[63:0]};
  \node[acc, below=of a, minimum width=22mm] (add) {add / sub};
  \node[acc, right=of add, minimum width=22mm] (log) {and / or / xor / not};
  \node[acc, right=of log, minimum width=22mm] (sh) {shifts};
  \node[acc, below=of log] (cmp) {compares};
  \node[trap, right=of cmp] (ov) {overflow};
  \node[trap, below=of ov] (ve) {value error};
  \draw[arr] (a) -- (add);
  \draw[arr] (b) -- (log);
  \draw[arr] (add) -- (cmp);
  \draw[arr] (sh) -- (ov);
  \draw[arr] (sh) -- (ve);
\end{tikzpicture}
\caption{Integer ALU (\texttt{pycore\_int\_alu}), combinational in the issue cycle.
It covers \texttt{+ - << >> \& | \textasciicircum\ \textasciitilde}, unary plus and minus, and the integer compares.
Beside the result it reports two flags.
\texttt{overflow\_flag} is set when add, subtract, unary minus, or a shift leaves the signed 64-bit range.
\texttt{value\_error} is a negative shift count.
Execute turns those into \texttt{PY\_TRAP\_OVERFLOW} and \texttt{PY\_TRAP\_VALUE}.
Compares reuse the subtractor and ignore the overflow flag.
A shift of 64 or more is the sign fill for \texttt{>>}, and either zero or overflow for \texttt{<<}, matching CPython's long-shift rules on a 64-bit limb.}
\label{fig:alu}
\end{figure}

\begin{figure}[htbp]
\centering
\begin{tikzpicture}[node distance=3.5mm]
  \node[rf] (b) {multiplier};
  \node[ctl, below=of b] (ch) {16-bit chunk\\top chunk signed};
  \node[acc, below=of ch] (pp) {65 $\times$ 17 partial product};
  \node[rf, left=8mm of pp] (a) {sext(multiplicand)};
  \node[acc, below=of pp] (acc) {acc $\leftarrow$ (acc $\gg$ 16) $+$ product};
  \node[rf, below=of acc] (lo) {low word, shifted in};
  \node[ctl, right=8mm of acc] (n) {4 steps $+$ register};
  \draw[arr] (b) -- (ch);
  \draw[arr] (ch) -- (pp);
  \draw[arr] (a) -- (pp);
  \draw[arr] (pp) -- (acc);
  \draw[arr] (acc) -- (lo);
  \draw[arr] (n) -- (acc);
\end{tikzpicture}
\caption{Signed 64$\times$64 multiplier (\texttt{pycore\_mul}, \texttt{STEP}$=16$).
Each cycle multiplies the full multiplicand by one chunk of the multiplier and folds it into a shift-right accumulator.
The multiplier is treated as unsigned chunks with only the top chunk signed, so the running sum is the exact signed product with no final correction.
Four steps plus the registered result are 5 cycles.
The low 64 bits are the Python product when it fits; a high half that is not the sign extension of the low half is \texttt{PY\_TRAP\_OVERFLOW}.
The same multiplier is the engine behind integer power.
\texttt{STEP} is the area-versus-latency knob; the rest of the core does not depend on it.}
\label{fig:mul}
\end{figure}

\begin{figure}[htbp]
\centering
\begin{tikzpicture}[node distance=2.6mm]
  \node[ctl, minimum width=56mm] (c0) {accept \quad sign and magnitude};
  \node[ctl, minimum width=56mm, below=of c0] (c1) {leading-zero count \quad align \quad load};
  \node[acc, minimum width=56mm, below=of c1] (run) {radix-4 restoring steps};
  \node[acc, minimum width=56mm, below=of run] (fix) {Python floor fix-up};
  \node[rf, minimum width=56mm, below=of fix] (out) {quotient and remainder};
  \node[trap, right=6mm of c0] (z) {div 0};
  \foreach \p/\q in {c0/c1,c1/run,run/fix,fix/out} \draw[arr] (\p) -- (\q);
  \draw[darr] (c0) -- (z);
  \node[anno, align=left, anchor=north west] at ($(run.east)+(0.2,-0.15)$)
    {three parallel\\subtractions\\of $d$, $2d$, $3d$\\two quotient bits\\per cycle};
\end{tikzpicture}
\caption{Floor divider (\texttt{pycore\_div}, \texttt{RL}$=2$) on \texttt{pycore\_udiv\_seq}.
Only $n = \mathrm{bits}(|a|) - \mathrm{bits}(|b|) + 1$ quotient bits are iterated, rounded up to a multiple of two.
Latency is $5 + \lceil n/2 \rceil$: 5 cycles when $|a| < |b|$, 37 for a full 64-bit quotient.
The five fixed cycles are accept, leading-zero and alignment, core load, core result, and the floor fix-up.
Division by zero is combinational in the accept cycle and does not start the unit.
\texttt{INT64\_MIN // -1} overflows.
The same restoring core with \texttt{SQRT}$=1$ is the floating-point square root.}
\label{fig:div}
\end{figure}

\begin{figure}[htbp]
\centering
\begin{tikzpicture}[node distance=3mm and 5mm]
  \node[rf] (e) {exponent};
  \node[ctl, below=of e] (z) {$e = 0$ \quad result 1};
  \node[ctl, right=of z] (bits) {walk bits MSB first};
  \node[acc, below=of bits] (sq) {square};
  \node[acc, right=of sq] (mu) {multiply if bit set};
  \node[acc, below=of sq] (mul) {shared \texttt{pycore\_mul}};
  \node[trap, right=of mul] (ov) {overflow};
  \draw[arr] (e) -- (z);
  \draw[arr] (z) -- (bits);
  \draw[arr] (bits) -- (sq);
  \draw[arr] (sq) -- (mu);
  \draw[arr] (sq) -- (mul);
  \draw[arr] (mu) -- (mul);
  \draw[arr] (mul) -- (ov);
  \node[anno, anchor=west] at ($(ov.east)+(0.2,0)$) {abort mid-chain};
\end{tikzpicture}
\caption{Integer power (\texttt{pycore\_ipow}).
A non-negative exponent does $\mathrm{bits}(e)-1$ squarings and $\mathrm{popcount}(e)-1$ extra multiplies, each a full multiplier pass, plus four cycles of sequencing.
A zero exponent finishes in the issue cycle.
A negative exponent is a float in Python and is re-routed before this unit.
The unit traps as soon as an intermediate product leaves the signed 64-bit range, so the long chains that complete are bases $0$, $1$, and $-1$.}
\label{fig:ipow}
\end{figure}
